% ============================================================
%  Chapter 5 — Validation
%  STATUS: §5.1 first draft (2026-05-27). §5.2..§5.9 are stub
%  sections so cross-references resolve; bodies to be drafted.
%  Source: CHAPTER5_BRIEF.md, DD-V5-032/033/034, evidence reports
%  under nexus-mvv-v5/results/.
% ============================================================
\chapter{Validation}
\label{ch:results}

This chapter validates the \NEXUS{} architecture of
Chapter~\ref{ch:architecture}, using the empirical evidence produced
by the implementation of Chapter~\ref{ch:implementation}. The work
done in this chapter is \emph{architectural validation}: each section
takes one architectural property promised in Chapter~\ref{ch:architecture}
--- bounded disclosure, joint welfare maximisation, the
$\sigma$-tier audit signature, and so on --- and asks whether the
property survives contact with the running system. The chapter is
not a contest in which \NEXUS{} is the contender and the baseline
mechanisms are opponents; the baselines are characterisation tools
for the architecture, and the relevant verdicts are about the
architecture's properties, not about which mechanism wins on
throughput.

The empirical surface is a single, deterministic, pre-registered
benchmark pipeline. It consists of $720$ simulation runs --- three
runtime mechanisms across eight benchmark cells, thirty seeds per
cell --- together with four specialised follow-on analyses that
unpack the audit log along axes (bits of disclosure, decision
attribution, centralised alignment, coordination-event latency) the
sweep alone does not surface. The chapter opens with the methodology
that produced this evidence; the seven sections that follow each
validate one architectural property against one or more pieces of
that evidence; the closing section frames the limits of what the
benchmark surface can validate.

% ------------------------------------------------------------
\section{Methodology}
\label{sec:val-methodology}

This section fixes the methodology for the rest of the chapter. It
states the framing under which validation is performed
(Section~\ref{subsec:val-framing}), enumerates the eight architectural
properties the chapter validates (Section~\ref{subsec:val-properties}),
introduces the five evidence families that supply the data
(Section~\ref{subsec:val-evidence}), names the mechanism configurations
under comparison (Section~\ref{subsec:val-mechanisms}), specifies the
statistical methods and thresholds (Section~\ref{subsec:val-stats}),
and gives the structural reading guide for the section sequence that
follows (Section~\ref{subsec:val-roadmap}).

\subsection{Framing: architectural validation, not outcome dominance}
\label{subsec:val-framing}

The thesis contribution stated in Chapter~\ref{ch:architecture} is a
\emph{reference architecture}: a layered composition of BDI Worker
agents, a multi-issue Marketplace mediator, an audit-only Fence and a
deferred School slot, together with the architectural properties that
composition is designed to deliver. The contribution is the
architecture and its properties, not a claim that this particular
implementation of the architecture wins on plant-floor outcomes
against every alternative on every benchmark. The distinction matters
for how validation is conducted, and so this chapter is structured
around it.

\emph{Architectural validation} asks: does the system, when run on
representative benchmarks, exhibit each of the architectural properties
promised by Chapter~\ref{ch:architecture}? Does the Marketplace
mediator actually find joint welfare maxima under an admissible
heuristic? Does the Worker layer's mono-dimensional-output discipline
translate into a measurable attribution pattern in the audit log?
Does the Fence layer produce a distinguishable safety signature
across mechanisms? These are property-level questions; their answers
characterise the architecture. \emph{Outcome dominance}, in contrast,
asks whether \NEXUS{} produces strictly better defect counts or
breakdown rates than a chosen baseline on a chosen benchmark family;
this is a question about one particular implementation on one
particular workload, not about the architecture as such.

The architectural-validation framing was not a retrospective
rephrasing of what the empirical work had already produced. The
contributions claimed in Chapter~\ref{ch:architecture} are stated as
architectural properties; an earlier pre-registered hypothesis surface
(an outcome-dominance test of \NEXUS{} against Siloed-v5, and a
productivity-differential test on centralised alignment) was
methodological scaffolding, ratified inside the sweep driver before
the data were inspected, and is treated in this chapter as
\emph{outcome characterisation}: the corresponding distributions and
verdicts are reported honestly in Section~\ref{sec:val-boundary} and
they do not bear on the architecture's claims (\texttt{DD-033},
\texttt{DD-034}). The other two pre-registered hypotheses --- the
alignment test and the disclosure test --- do bear on architectural
properties and are load-bearing for
Sections~\ref{sec:val-alignment} and~\ref{sec:val-disclosure}
respectively. Section~\ref{subsec:val-stats} names all of them
explicitly, fixes the descriptive labels used throughout the
chapter, and states what would falsify each.

\subsection{The eight architectural properties under validation}
\label{subsec:val-properties}

Chapter~\ref{ch:architecture} promises eight architectural properties,
each of which is the subject of one of this chapter's later sections.
Some are stated explicitly in the architectural prose; others are
structural commitments the prose makes implicitly, in the sense that
the architecture cannot be implemented without them holding. We
enumerate them here in the order in which the chapter validates them,
not in their order of appearance in Chapter~\ref{ch:architecture}.
Table~\ref{tab:val-properties} summarises the mapping from each
property to the chapter section where it is promised and to the
section of this chapter that validates it.

\begin{description}[leftmargin=2.4em,style=nextline,itemsep=0.25em]
  \item[P1 --- Layered separation of concerns.] The four ISA-95
    Level~3 functions are decomposed into four Worker agents, and
    coordination, governance, and adaptation are separate layers
    that can evolve independently. P1 is foundational --- it is not
    validated in isolation, but every result in
    Sections~\ref{sec:val-emergence}--\ref{sec:val-attribution}
    is conditional on the layer boundaries holding.
  \item[P2 --- Mono-dimensional outputs and domain authority.] Each
    Worker emits constraints primarily on its home issue, with
    cross-domain constraints damped to act as tie-breakers. Validated
    in Section~\ref{sec:val-attribution} by the decision-attribution
    analysis (the $4\times 4$ contribution matrix).
  \item[P3 --- Bounded disclosure.] Agents publish only bounded bid
    summaries; never beliefs, constraints, or internal weights.
    Validated in Section~\ref{sec:val-disclosure} by the bits-per-event
    disclosure quantification across mechanisms.
  \item[P4 --- Joint welfare maximisation via admissible search.]
    The mediator's branch-cut A* with an admissible heuristic
    (\texttt{DD-014}) selects the contract that maximises the sum of
    agent utilities over the bid space. Validated in
    Section~\ref{sec:val-alignment} by the alignment of \NEXUS{}'s
    chosen contracts against the centralised utility-sum ceiling
    (CUSO, \texttt{DD-032}).
  \item[P5 --- $\sigma$-tier governance signature.] The Fence
    classifies each contract into one of three tiers
    ($\sigma_0$/$\sigma_1$/$\sigma_2$); coordination architectures
    produce distinguishable tier profiles. Validated in
    Section~\ref{sec:val-sigma} by the cross-mechanism $\sigma$-tier
    counts produced by the main coordination sweep.
  \item[P6 --- Coordination over independence.] \NEXUS{} produces
    contracts that are characteristically different from those
    produced by an uncoordinated, home-only baseline running on the
    same code and the same environment. Validated \emph{cumulatively}
    by the $\sigma$-tier signature (Section~\ref{sec:val-sigma}), the
    disclosure quantification (Section~\ref{sec:val-disclosure}), the
    centralised alignment (Section~\ref{sec:val-alignment}) and the
    decision-attribution analysis (Section~\ref{sec:val-attribution});
    P6 is not validated by a single outcome contrast (the relevant
    outcome characterisation is in Section~\ref{sec:val-boundary}).
  \item[P7 --- Compositional emergence.] Joint trade-offs the agents
    themselves never explicitly programme --- a rush order against a
    constrained maintenance budget yielding a deferred maintenance
    window, a premium product against a borderline defect rate
    yielding a partial quality hold --- emerge from the mediator's
    composition. Validated in Section~\ref{sec:val-emergence}.
  \item[P8 --- Order-invariant decisions.] The mediation outcome
    does not depend on the order in which the agents speak, because
    the mediator operates on the complete bid set in a single round.
    Validated in Section~\ref{sec:val-order} by the agent-order
    permutation test.
\end{description}

A ninth architectural slot --- the School layer, the locus of
continual preference adaptation --- is named in
Chapter~\ref{ch:architecture} as future work and is not validated
empirically here; it is taken up in
Chapter~\ref{ch:conclusions}.

\begin{table}[ht]
\centering
\caption{The eight architectural properties under validation, the
section of Chapter~\ref{ch:architecture} that promises each, and the
section of this chapter that validates each.}
\label{tab:val-properties}
\small
\begin{tabularx}{\textwidth}{@{}l X l l@{}}
\toprule
\textbf{Property} & \textbf{Statement} & \textbf{Promised} & \textbf{Validated} \\
\midrule
P1 & Layered separation of concerns
   & Sec.~\ref{sec:arch-overview}, \ref{subsec:why-layered}
   & \S\ref{sec:val-emergence}--\ref{sec:val-attribution} (implicit) \\
P2 & Mono-dimensional outputs + domain authority
   & Sec.~\ref{subsec:mono-dim}
   & \S\ref{sec:val-attribution} \\
P3 & Bounded disclosure
   & Sec.~\ref{subsec:bids}, \ref{subsec:privacy}
   & \S\ref{sec:val-disclosure} \\
P4 & Joint welfare maximisation via admissible search
   & Sec.~\ref{subsec:mediator}
   & \S\ref{sec:val-alignment} \\
P5 & $\sigma$-tier governance signature
   & Sec.~\ref{subsec:fence-tiers}, \ref{subsec:fence-logonly}
   & \S\ref{sec:val-sigma} \\
P6 & Coordination over uncoordinated baseline
   & Sec.~\ref{sec:arch-intro}, \ref{subsec:protocol-scope}
   & \S\ref{sec:val-sigma}--\ref{sec:val-attribution} (cumulative) \\
P7 & Compositional emergence
   & Sec.~\ref{subsec:mediator} + Sec.~\ref{sec:impl-baselines}
   & \S\ref{sec:val-emergence} \\
P8 & Order-invariant decisions
   & Sec.~\ref{subsec:protocol-scope}
   & \S\ref{sec:val-order} \\
\bottomrule
\end{tabularx}
\end{table}

\subsection{Evidence families}
\label{subsec:val-evidence}

The empirical work draws on five distinct \emph{evidence families}.
Each family supplies a particular slice of evidence, anchored in a
particular property of the architecture and analysed by a dedicated
post-processor. Figure~\ref{fig:val-stack} shows the architectural
layer each family principally bears on, together with the property
labels.

\begin{description}[leftmargin=2.4em,style=nextline,itemsep=0.25em]
  \item[Main coordination sweep.] The workhorse evidence family. Three
    mechanisms (\NEXUS{}, Siloed and Rule-Based Coordination) are run
    across eight benchmark cells (two Brandimarte instances --- Mk06
    and Mk10 --- and six \texttt{synthetic\_v2} cells covering the
    \{\texttt{medium}, \texttt{large}\} shop sizes crossed with the
    \{\texttt{low}, \texttt{medium}, \texttt{high}\} coupling levels)
    at $30$ paired seeds per (mechanism, cell), for a total of $720$
    runs. Each run yields outcome metrics, $\sigma$-tier classifier
    output and the full Contract Record audit log. The sweep is the
    evidence base for Sections~\ref{sec:val-emergence},
    \ref{sec:val-order}, \ref{sec:val-sigma} and~\ref{sec:val-attribution}.
  \item[Decision attribution analysis.] A specialised analysis of the
    Contract Records that computes, per mechanism and per cell, a
    $4\times 4$ contribution matrix recording how often each agent's
    highest-utility bid pins each of the four contract dimensions in
    the final composite contract. The diagonal mass measures
    home-issue authority share; the off-diagonal mass measures
    cross-domain influence. Evidence base for
    Section~\ref{sec:val-attribution} (property P2).
  \item[Centralised alignment benchmark.] A post-hoc enumeration of
    the entire $108$-contract joint space for every conflict in the
    sweep, evaluating the centralised utility-sum optimum CUSO
    (\texttt{DD-032}). For each conflict the per-conflict alignment
    ratio \NEXUS{}\,/\,\textsc{cuso} measures the fraction of the
    centralised utility-sum ceiling that the bounded-disclosure
    distributed mechanism recovers. Evidence base for
    Section~\ref{sec:val-alignment} (property P4).
  \item[Disclosure quantification.] A bit-budget calculation applied
    to every bid and to the CUSO reference, instantiating the
    information-disclosure model of \texttt{DD-030}. Reported per
    mechanism, per agent, and as the per-event ratio of mechanism
    disclosure to CUSO disclosure. Evidence base for
    Section~\ref{sec:val-disclosure} (property P3).
  \item[Coordination-event latency.] A wall-time measurement per
    coordination event, used in Section~\ref{sec:val-boundary} as a
    deployment-feasibility note against the ISA-95 Level~3 control
    cycle. It is not a property validation.
\end{description}

\nexustikzfigure{fig:val-stack}{The four-layer \NEXUS{} architecture
alongside the five evidence families that populate this chapter. The
property tags shown on each side make the validation correspondence
explicit; Section~\ref{subsec:val-roadmap} sequences the sections
that exploit the correspondence.}{fig_val_stack}

The thesis-facing names used above and throughout the chapter
intentionally differ from the internal working names that appear in
the repository (Step~1 sweep, Phase~C/D/E/F analyses). The renaming
makes the chapter readable as a self-contained empirical narrative
without the bookkeeping baggage of the engineering process that
produced the data.

\subsection{Mechanism configurations evaluated}
\label{subsec:val-mechanisms}

The architecture is characterised against three drop-in alternatives
to its Marketplace layer. Each alternative isolates a single
architectural axis, so that the contrasts of
Sections~\ref{sec:val-emergence}--\ref{sec:val-attribution} are
single-variable contrasts and the differences observed are
attributable to the architectural property under test.
Table~\ref{tab:val-mechanisms} places the four mechanisms on the two
axes that the comparison turns on: whether the mechanism performs a
joint search over the contract space, and whether it has access to
the full state of every agent's utility function.

\begin{description}[leftmargin=2.4em,style=nextline,itemsep=0.25em]
  \item[\NEXUS{}.] The architecture under validation
    (Chapter~\ref{ch:architecture}, Chapter~\ref{ch:implementation}).
    Joint search over the published bid sets; bounded disclosure.
  \item[Siloed-v5.] The uncoordinated baseline of
    Section~\ref{subsec:impl-siloed}. Each agent emits home-only bids
    (\texttt{cross\_domain\_awareness=False}) and the mechanism
    composes their independent home picks into a single contract.
    No joint search; bounded disclosure; same code and same
    environment as \NEXUS{}, so the contrast with \NEXUS{} is a
    single-variable contrast on coordination.
  \item[RBC --- Rule-Based Coordination.] The
    static-threshold-rule baseline of \texttt{DD-021} and
    Section~\ref{subsec:impl-rbc}. Each contract dimension is
    decided from a fixed per-dimension threshold on the conflict
    announcement; no agents, no bids. RBC trades coordination
    flexibility for determinism and a known rule table.
  \item[CUSO --- Centralised Utility-Sum Optimiser.] The
    full-information ceiling of \texttt{DD-032} and
    Section~\ref{subsec:impl-cuso}. CUSO enumerates the entire
    $108$-contract joint space and returns the contract that
    maximises the sum of agent utilities, with access to every
    agent's full utility function. It is not a runtime mechanism;
    it is a post-hoc reference applied to each conflict in the
    sweep.
\end{description}

The Sequential cascade of \texttt{DD-012} and the single-winner
Contract Net of \texttt{DD-001} remain in the code but are not used
as primary baselines: the Sequential cascade was shown to exhibit a
structural commitment-asymmetry under resource-constrained loads
(\texttt{DD-019}) that confounds the comparison on outcomes the
thesis cares about, and the single-winner Contract Net's output is a
winner-label string that cannot, by construction, represent a
cross-functional trade-off (\texttt{DD-001}).

\begin{table}[ht]
\centering
\caption{The four mechanism configurations under comparison, placed
on the joint-search and information-access axes that the contrasts of
this chapter turn on.}
\label{tab:val-mechanisms}
\small
\begin{tabularx}{\textwidth}{@{}l c c X@{}}
\toprule
\textbf{Mechanism} & \textbf{Joint search} & \textbf{Information access} & \textbf{Role in the comparison} \\
\midrule
\NEXUS{}    & yes (branch-cut A*) & bounded (bid summaries)
            & architecture under validation \\
\addlinespace
Siloed-v5   & no (home-only picks) & bounded (bid summaries)
            & isolates coordination vs independence \\
\addlinespace
RBC         & no (static rule table) & public rule table
            & isolates dynamic vs static coordination \\
\addlinespace
CUSO        & full ($108$-contract enumeration) & full agent utility functions
            & centralised reference ceiling \\
\bottomrule
\end{tabularx}
\end{table}

\subsection{Statistical methods and thresholds}
\label{subsec:val-stats}

Outcome distributions across the $30$ paired seeds per cell are
compared with the two-sided \emph{Mann-Whitney U} test --- a
non-parametric rank test that makes no distributional assumption and
that is appropriate to the small, paired and skewed outcome
distributions that the discrete-event simulator produces. Effect size
is reported as \emph{Cliff's $\delta$}, which takes values in
$[-1, +1]$ and admits the conventional thresholds: $|\delta|\,\ge\,0.147$
small, $0.330$ medium and $0.474$ \emph{substantively large}. A positive
$\delta$ in a contrast ``A vs B'' means that A's outcome distribution
is shifted in the favoured direction relative to B's; the favoured
direction is specified per outcome (lower is better for defects,
breakdowns and makespan; higher is better for throughput). The
significance level is Bonferroni-corrected against the multiple
comparisons induced by four outcomes and two pairwise contrasts:
$\alpha = 0.05/8 = 0.00625$.

The empirical surface of the chapter is organised around eight
named tests, each anchored to one of the architectural properties
P1--P8 of Section~\ref{subsec:val-properties}. Four of the tests
were pre-registered in
\texttt{experiments/run\_step1\_unified\_sweep.py} before any sweep
data were inspected; the other four are structural or
regression-anchored checks. The chapter refers to all of them by
descriptive two-letter labels so the prose stays legible without
the bookkeeping of the internal hypothesis codes carried by the
code, the JSON keys of the result files, and the design-decision
ledger. The mapping is one-to-one and is recorded in
Table~\ref{tab:val-tests} below; the labels are the chapter's
glossary and are used uniformly throughout
Sections~\ref{sec:val-emergence}--\ref{sec:val-boundary}.

\begin{table}[ht]
\centering
\caption{The eight tests under which the chapter's empirical surface
is reported. The label is the chapter glossary; the internal name is
the identifier used by the code, the result-file JSON, and the
design-decision ledger; ``Load-bearing'' marks the tests whose
verdict bears on an architectural claim, as opposed to outcome
characterisations whose verdicts are reported in
Section~\ref{sec:val-boundary} per \texttt{DD-033}/\texttt{DD-034}.}
\label{tab:val-tests}
\small
\begin{tabularx}{\textwidth}{@{}l l X c c c@{}}
\toprule
\textbf{Label} & \textbf{Test name} & \textbf{What is tested}
  & \textbf{Section} & \textbf{Pre-reg.} & \textbf{Load-bearing} \\
\midrule
OD   & Outcome dominance
     & \NEXUS{} outcome distributions strictly dominate Siloed-v5 on
       quality outcomes
     & \S\ref{sec:val-boundary}
     & yes (HB1)
     & char.\ \\
AL   & Alignment
     & \NEXUS{}'s per-conflict alignment against the CUSO ceiling has
       $p_{25} \ge 0.945$ on every cell
     & \S\ref{sec:val-alignment}
     & yes (HF1)
     & yes \\
DI   & Disclosure
     & \NEXUS{}'s per-event disclosure is at most $30\%$ of the
       CUSO ceiling on every cell
     & \S\ref{sec:val-disclosure}
     & yes (HF2)
     & yes \\
PR   & Productivity differential
     & \NEXUS{} exceeds Siloed-v5 in CUSO alignment by at least
       $5$ percentage points on every cell
     & \S\ref{sec:val-boundary}
     & yes (HF3)
     & char.\ \\
$\sigma$G & $\sigma$-tier governance
     & The three runtime mechanisms produce substantively ordered
       $\sigma$-tier counts (not indistinguishable noise)
     & \S\ref{sec:val-sigma}
     & yes (HA1)
     & yes \\
OI   & Order invariance
     & \NEXUS{} produces byte-identical contracts under any
       permutation of agent presentation order
     & \S\ref{sec:val-order}
     & no (structural)
     & yes \\
CE   & Compositional emergence
     & F6 + F7 trade-offs emerge under \NEXUS{} mediation, absent
       under Siloed-v5 and RBC
     & \S\ref{sec:val-emergence}
     & no (regression)
     & yes \\
JM   & Joint mediation signature
     & \NEXUS{}'s diagonal contribution mass differs from Siloed-v5's
       by at least $10$ percentage points uniform across cells
     & \S\ref{sec:val-attribution}
     & no (structural)
     & yes \\
\bottomrule
\end{tabularx}
\end{table}

Two of the four pre-registered tests --- OD (the outcome-dominance
test) and PR (the productivity-differential test) --- are explicitly
reclassified as outcome characterisation per \texttt{DD-033} and
\texttt{DD-034}: their verdicts are reported in
Section~\ref{sec:val-boundary}, but they do not bear on the
architecture's claims, because the architecture of
Chapter~\ref{ch:architecture} does not promise outcome dominance over
a same-codebase uncoordinated baseline. The remaining six tests bear
on the architecture's property surface and are load-bearing for the
section in which they run. The non-pre-registered tests (OI, CE, JM)
are not statistical-distribution tests: OI is a structural property
checked by a determinism regression, CE is a qualitative-emergence
inference anchored by the \texttt{DD-023} regression suite, and JM
is an architecture-level $10$-percentage-point reading on the
decision-attribution matrix. Each is introduced in the section that
runs it.

\subsection{Reading guide for the rest of the chapter}
\label{subsec:val-roadmap}

The remainder of the chapter follows a uniform structure. Each of
Sections~\ref{sec:val-emergence}--\ref{sec:val-attribution} takes
one architectural property, names the evidence family it draws on,
states the measurement performed, reports the empirical result, and
closes with a paragraph on what the result establishes and what it
does not. The order of the sections is from the compositional and
order-invariance claims that follow from the joint protocol
(Sections~\ref{sec:val-emergence} and~\ref{sec:val-order}), through
the governance signature (Section~\ref{sec:val-sigma}), the two
quantitative architectural claims of bounded disclosure and
centralised alignment (Sections~\ref{sec:val-disclosure}
and~\ref{sec:val-alignment}), to the decision-attribution signature
of the mediation itself (Section~\ref{sec:val-attribution}). The
six sections collectively exhaust the architectural surface
promised by Chapter~\ref{ch:architecture}.

Property P6 (coordination over independence) deserves a separate
note here. The architecture's claim that joint mediation produces
contracts \emph{recognisably different} from those produced by an
uncoordinated baseline is supported \emph{cumulatively} ---
distinct $\sigma$-tier signature (Section~\ref{sec:val-sigma}),
distinct disclosure profile (Section~\ref{sec:val-disclosure}),
distinct alignment ratios against the centralised ceiling
(Section~\ref{sec:val-alignment}), distinct decision-attribution
mass (Section~\ref{sec:val-attribution}). It is not validated by a
single outcome contrast against Siloed-v5; on the FJSP benchmark
families the within-codebase outcome distributions are at parity,
and that result is part of the outcome characterisation honestly
reported in Section~\ref{sec:val-boundary}. The four cumulative
contrasts taken together are what makes the property load-bearing.

Section~\ref{sec:val-boundary} closes the chapter with the
applicability boundary. There are two reasons the section is given
its own space rather than woven into the earlier sections. The first
is that the outcome characterisation honestly belongs together with
the limits of what benchmark validation can prove --- the
reclassified outcome-dominance and productivity verdicts are not
isolated empirical observations but features of a study that is, by
design, about the architecture rather than about which specific
implementation of the architecture wins on which specific scalar
metric. The second is that the chapter's claims
need to be stated against the design space they implicitly occupy.
The \NEXUS{} architecture admits a range of implementation choices on
each layer --- BDI workers could be replaced by a multi-agent
reinforcement-learning component; the multi-issue A* mediator could
be replaced by a combinatorial auction; the log-only Fence could be
replaced by an enforcement layer with runtime intervention --- and the
architectural properties validated here are properties of the
\emph{reference} architecture, not of any one specific implementation.
Section~\ref{sec:val-boundary} makes that design space explicit, names
the limitations of the present study against it, and previews how
Chapter~\ref{ch:industrial}'s industrial case study addresses the
gap between benchmark validation and operational deployment. The
chapter does not, anywhere, claim outcome dominance over a baseline
on any benchmark; that the architecture can be characterised against
its baselines without such a claim being made is itself a property
of the chapter's framing, and Section~\ref{sec:val-boundary} is where
that framing is closed out.

% ============================================================
%  §5.2 .. §5.9  --- stubs (labels only; bodies to be drafted)
% ============================================================

\section{Compositional trade-off emergence}
\label{sec:val-emergence}

The second architectural property under validation is P7:
cross-functional trade-offs emerge from the Marketplace mediator's
joint composition of the agents' bids, rather than from explicit
trade-off rules programmed into any single component. P7 is a
structural claim about \emph{where} cross-functional decisions are
located. If a useful trade-off pattern --- ``defer maintenance when
the budget is tight \emph{and} the job is a rush order''; ``hold
quality more strictly when the product variant is more expensive''
--- has to be coded somewhere as a special-case rule, then the
mediator is not adding value; it is just executing rules its
designers anticipated. P7 says the trade-offs emerge from the A*
composition of independently weighted home-domain constraints,
without any agent or rule table writing the trade-off down.

\subsection*{Evidence and inference pattern}

This is qualitative-emergence evidence, not a statistical
distribution test, and the pattern of inference is structural: a
trade-off is present where mediation is present and absent where
mediation is absent. Two concrete trade-offs are anchored by the
\texttt{DD-023} regression suite
(\texttt{tests/test\_dd\_v5\_023\_extensions.py}, $40$ tests):

\begin{description}[leftmargin=2.4em,style=nextline,itemsep=0.25em]
  \item[F6 --- rush job under a tight maintenance budget defers
    maintenance.] A conflict where the job is a rush job
    (\texttt{is\_rush}~$=$~True) on a degraded machine
    (\texttt{machine\_health}~$=$~$0.45$), and the maintenance budget is
    nearly exhausted (\texttt{remaining\_budget\_ratio}~$=$~$0.1$),
    yields a contract whose maintenance window is \textsc{planned},
    \emph{not} \textsc{immediate} --- even though the machine's
    health alone would normally argue for \textsc{immediate}, and
    even though the rush job alone would normally tighten the
    decision in the other direction. The trade-off is the
    composition of the two pulls under the budget constraint.
  \item[F7 --- premium product on a borderline defect rate
    triggers a partial quality hold.] A conflict where the product
    variant is \texttt{premium} (high unit value) and the defect
    rate sits just below the quality threshold ($0.04$) yields a
    contract whose quality action is \textsc{hold\_light}, rather
    than the \textsc{no\_hold} that the defect rate alone would
    select.
\end{description}

The regression suite pins these compositions as targeted unit-style
tests: each test constructs a \texttt{ConflictAnnouncement} with the
relevant joint signal, runs the agents' bid generation against the
mediator, and asserts that the resulting weight ordering or chosen
contract value reflects the trade-off. The tests are not the
discovery --- they are the regression guard for the discovery, and
the discovery itself is in the design of the agents' constraint
catalogues plus the A* mediator's composition rule (the priority key
$f = g + h$ of Section~\ref{subsec:val-stats} applied across all four
agents' bid utilities). The trade-off emerges from the structural
arithmetic of the joint search; the test is what tells us it is still
emerging after a refactor.

The F6 mechanism is worth tracing in detail because the structural
arithmetic is what the chapter is asking the reader to grant;
Figure~\ref{fig:val-emergence} shows the path.
\PMA{}'s
constraint catalogue contains two relevant unary constraints on the
maintenance dimension: \texttt{pma\_force\_IMMEDIATE\_when\_critical}
and \texttt{pma\_prefer\_PLANNED\_when\_degraded}. The
industrial-awareness extensions of \texttt{DD-023} add two
modulators: the priority multiplier (rush job $\to$ multiplicative
boost on both constraints, weighting maintenance up overall) and the
budget discount (low \texttt{remaining\_budget\_ratio} $\to$
multiplicative discount on the \textsc{immediate} constraint only,
because it is the costly action). On the F6 signal (rush,
\texttt{machine\_health} $= 0.45$, \texttt{remaining\_budget\_ratio} $=
0.1$), the priority multiplier scales both constraints by $5\times$;
the budget discount factor of about $0.2$ applies only to
\textsc{immediate}; and the resulting weight ordering inverts:
\texttt{pma\_prefer\_PLANNED\_when\_degraded} now carries the heavier
weight in \PMA{}'s bid set, and the mediator's joint search picks the
bid combination whose joint region admits \textsc{planned}. No code
path anywhere computes ``if rush and budget tight, then PLANNED''.
The weight modulators are independent; the budget discount lives on
the \textsc{immediate} constraint and is unaware of the priority
multiplier; the joint outcome is the sum the mediator finds across
the four agents' bid utilities. F7 has the same structural shape:
the multi-product modulator weights \QCA{}'s
\texttt{qca\_partial\_hold\_when\_borderline} constraint conditional
on \texttt{product\_unit\_value}, and the mediator's joint composition
chooses \textsc{hold\_light} where the base rule would have chosen
\textsc{no\_hold}.

\nexustikzfigure{fig:val-emergence}{The F6 compositional-emergence
path. A rush conflict on a degraded machine under a near-exhausted
maintenance budget inverts the weight ordering of \PMA{}'s two
maintenance constraints --- the priority multiplier boosts both, the
budget discount suppresses only \textsc{immediate} --- so the
mediator's joint search admits \textsc{planned} maintenance. No
component encodes the trade-off explicitly; it is the arithmetic of
the composition.}{fig_val_emergence}

\subsection*{Absence under Siloed and RBC}

The same trade-offs do not emerge under Siloed-v5. Siloed-v5
constructs the contract by composing each agent's independent
home-issue pick: \PMA{} chooses the maintenance window from
\texttt{machine\_health} alone, \QCA{} chooses the quality action
from \texttt{defect\_rate} alone, and so on. The rush flag and the
budget ratio reach \PMA{}; they do not, by construction, reach
\QCA{}'s home decision; and \PSA{}'s production rate has no way to
influence the maintenance choice. There is no place in Siloed-v5
where cross-domain trade-offs can compose, so F6 and F7 do not
emerge there. This is not a Siloed-v5 bug; it is what ``no joint
mediation'' means.

The same holds for RBC, with one nuance. RBC's threshold rule table
(Section~\ref{subsec:impl-rbc}) explicitly encodes a small handful
of priority-and-budget interactions on a per-dimension basis ---
the \texttt{TestRBCRuleOrderDependence} regression test fixes the
correct application order (priority first, then budget) as a
permanent guard against silent rule-rewrite. RBC therefore exhibits
\emph{programmed} trade-offs on the dimensions its rule table was
extended to anticipate. RBC's trade-offs are not emergent; they are
the rule designer's explicit decisions encoded as conditionals. P7
asks for trade-offs that emerge \emph{without} being coded as such,
and that property is supported only on \NEXUS{}: the same F6 and F7
patterns arise from the A* composition without any rule in any
agent or in any rule table that says ``if rush and budget tight,
then PLANNED''.

\subsection*{Architecture-level proxy on the main sweep}

The \texttt{DD-023} trade-offs are pinned on the agent constraint
catalogues by the regression suite, but the \texttt{synthetic\_v2}
benchmark cells of the main sweep do not exercise the industrial
signals (\texttt{remaining\_budget\_ratio}, \texttt{product\_type})
that F6 and F7 turn on --- those fields are populated only on the
industrial-automotive environment path of
Chapter~\ref{ch:industrial}. The compositional-emergence signature
on the main sweep is therefore architecture-level rather than
trade-off-specific: it is the $10$~percentage-point reduction in
\NEXUS{}'s diagonal contribution mass relative to Siloed-v5's that
Section~\ref{sec:val-attribution} reports. The diagonal mass measures
the share of the final contract whose value coincides with the
chosen bid's \emph{home-issue} sample; a reduction of that share
under joint mediation is what compositional emergence looks like at
the contribution-attribution level, because the mediator is pulling
the chosen bids' \emph{non-home} dimensions into the joint region.

The uniformity of the $10$-pp gap across all $8$ cells is itself
the property-level evidence for P7 on the main sweep. The
trade-off-specific patterns F6 and F7 are evidence on the
constraint-catalogue level; the contribution-attribution pattern is
the evidence at the architectural level.

\subsection*{Limitations}

Two are worth naming. First, F6 and F7 are two trade-offs ---
chosen because the design of the industrial agent extensions
(\texttt{DD-023}) explicitly anticipates them. The agent
constraint catalogues admit a larger family of compositions
(e.g., rush plus high defect rate plus low budget), and only a
subset has been pinned by the regression suite. Whether
\NEXUS{}'s mediator produces the architecturally desired
composition on every such joint signal is not exhaustively tested.

Second, the architecture-level proxy --- the $10$-pp diagonal-mass
gap --- is uniform across cells. Uniformity is informative
(P7 is a structural property of joint mediation, not a
benchmark-dependent accident) and also limiting (the proxy carries
no per-cell granularity that would let one cell argue more or less
strongly for emergence than another). The trade-off-specific
discussion lives in the regression suite; the
architecture-level discussion lives in
Section~\ref{sec:val-attribution}.

% ------------------------------------------------------------
\section{Order-invariant decisions}
\label{sec:val-order}

The third architectural property under validation in this batch is
P8: the contract \NEXUS{} produces on a given conflict does not
depend on the order in which the agents are presented to the
mediator. P8 is a structural property of the single-round protocol
of Section~\ref{subsec:protocol-scope} --- not an emergent property
that has to be discovered, but a consequence of the mediator's
algorithmic design that the implementation either preserves or
breaks.

\subsection*{The structural argument}

\NEXUS{}'s Marketplace, per Sections~\ref{subsec:mediator}
and~\ref{subsec:impl-mediator}, operates in a single round. Every
agent generates its complete bid set against the same
\texttt{ConflictAnnouncement}, the bid sets are collected on the
Bid Bus into a list-of-lists indexed by agent, and the mediator's
branch-cut A* search expands the combinatorial space of one-bid-per-agent
combinations under the priority key $f = g + h$. The bid-list order
is positional metadata: the search treats the lists as the input
to a combination enumeration, and the combination set
$\{(b_1, b_2, b_3, b_4) \mid b_i \in B_i\}$ is, by
definition, order-invariant in the agent index. Whichever agent
appears first in the list, the same set of combinations is
explored, and the priority key $f$ is the same scalar at every
combination because $g$ is a sum of bid utilities and $h$ is a sum
of per-agent suffix maxima --- both invariant under reordering of
the agents.

The tie-break inside the chosen combination's joint region is
deterministic (Section~\ref{subsec:impl-mediator}): the contract
that maximises the summed agent utility wins, with ties broken by
the canonical enumeration order of the $108$-contract space. The
canonical order is fixed at module import time and does not depend
on the agent presentation order. The full pipeline is therefore a
function of the bid sets and the announcement, with no order-of-input
dependence; the property P8 is structurally guaranteed and
verifying it empirically is a determinism check on the implementation,
not a discovery.

\subsection*{Empirical verification}

The implementation's determinism is verified by the test suite of
Section~\ref{sec:impl-testing}. \texttt{tests/test\_determinism.py}
exercises (a) cross-call determinism: two consecutive calls of
\texttt{agent.propose\_bids} on the same announcement produce
identical bid utilities and identical sample contracts; (b)
cross-process determinism: a fresh subprocess produces the same
per-agent random seeds, the same bid sets, and the same composed
contracts as the parent process. The cross-process test is the
critical guarantee, because it pins the architecture's
order-invariance against the specific implementation choice that
might break it --- per-process random-number-generator state. The
\texttt{\_stable\_seed} hash of
Section~\ref{subsec:impl-agentbase} replaces the standard library's
process-keyed hash and is what makes the regression test pass.

Combined with the structural argument, this gives a tight verdict on
P8 for \NEXUS{}: byte-identical contracts on the same announcement,
across processes, with the agent presentation order irrelevant to
the result. Figure~\ref{fig:val-order} contrasts the single-round
protocol that makes this hold with the sequential cascade that
breaks it.

\nexustikzfigure{fig:val-order}{Order invariance by contrast. In
\NEXUS{} (left) all four agents bid into the Bid Bus and the mediator
searches the complete bid set in a single round, so the composed
contract is byte-identical under any agent order. In the sequential
cascade (right, \texttt{DD-012}) each agent's pick restricts the next
agent's space, so a different order yields a different contract.}{fig_val_order}

\subsection*{Contrast with sequential cascade}

P8 is falsifiable: a mechanism in the same architecture but with a
different mediator can violate it. The Sequential cascade of
Section~\ref{subsec:impl-baselines} (\texttt{DD-012}) is the
documented example. The cascade fixes the agent order to
PSA $\to$ \PMA{} $\to$ \QCA{} $\to$ \SCA{}, and each agent's pick
restricts the contract space the next agent sees: \PSA{} commits
the production rate; \PMA{} chooses the maintenance window
conditioned on that production rate; \QCA{} chooses quality
conditioned on the first two; \SCA{} closes the loop. The
regression test
\texttt{tests/test\_baseline\_sequential.py::test\_cascade\_order\_psa\_pma\_qca\_sca}
pins this order even when the agent list is registered in reverse,
and the companion test
\texttt{test\_upstream\_fixed\_values\_constrain\_downstream}
verifies that an upstream pick can force a downstream agent into an
otherwise non-preferred choice (P\PSA{}'s rate fixes the joint
region, forcing \PMA{} into a maintenance window it would not
have selected first). The cascade therefore satisfies, by
construction, the negation of P8: the contract produced depends on
the agent order.

A different falsifying mechanism is RBC, on a smaller surface. The
\texttt{TestRBCRuleOrderDependence} regression in
\texttt{tests/test\_dd\_v5\_023\_extensions.py} pins the order in
which RBC applies its own internal rule adjustments (priority
first, then budget). A budget-first ordering would, on the test's
joint signal, downgrade the maintenance window differently and
produce a different contract. RBC therefore exhibits an
\emph{internal} rule-order dependence: the mechanism's authors had
to choose an order and the regression test makes that choice
permanent. P8 is satisfied by \NEXUS{} structurally; in RBC it is
satisfied only on the dimensions the rule table does not
sub-divide, and not by construction.

\subsection*{Limitation}

P8 is a property of the single-round protocol. \texttt{DD-006}
restricts \NEXUS{} to single-round mediation; the iterative-narrowing
extension of Hattori, Klein and Ito (Section~\ref{subsec:protocol-scope})
would break P8 by construction, because iterative rounds couple the
agents' bids to the partial-agreement feedback produced in earlier
rounds, and the order in which agents see and respond to that
feedback would influence the subsequent rounds. \NEXUS{}'s P8
guarantee therefore depends on the protocol-scope choice
\texttt{DD-006} and would be replaced by a softer guarantee (e.g.,
convergence-order invariance, or invariance up to a tie-break) under
iterative narrowing. Iterative narrowing is deferred as future work
in Chapter~\ref{ch:conclusions}; the single-round protocol is the
scope under which the chapter's claims are made, and P8 is one of
the chapter's strongest claims because the architectural argument and
the empirical regression test agree exactly on what they assert.

\section{$\sigma$-tier governance signature}
\label{sec:val-sigma}

The fourth architectural property under validation in this chapter is
P5: the Fence layer classifies every resolved contract into one of
three tiers ($\sigma_0$ safe, $\sigma_1$ at boundary, $\sigma_2$
violation), and different coordination mechanisms running on the same
plant under the same Fence rules produce \emph{distinguishable} tier
profiles. The property is what makes the Fence audit stream a
diagnostic surface for the architecture: if every mechanism produced
the same tier distribution, the Fence output would carry no
mechanism-specific information and could not be a useful audit signal.
The $\sigma$G test of Section~\ref{subsec:val-stats} is the
architectural-reading test that the three runtime mechanisms ---
\NEXUS{}, Siloed-v5 and RBC --- produce substantively ordered tier
profiles, not indistinguishable noise.

\subsection*{Evidence and measurement}

The evidence is the $\sigma$-tier classifier output emitted by the
Fence layer (Section~\ref{sec:impl-fence}) over the main coordination
sweep. For each of the $720$ runs the classifier evaluates every
resolved contract against the default four-rule safety suite of
Section~\ref{subsec:fence-tiers}
(\texttt{no\_fast\_when\_breakdown\_imminent},
\texttt{maintain\_when\_health\_critical},
\texttt{hold\_quality\_when\_severe},
\texttt{preorder\_when\_long\_lead}) and emits a
\texttt{FenceRecord} with the assigned tier, the rule names that
fired (violation and boundary), and the plant signals against which
the rules were evaluated. The aggregate of interest is the
\emph{$\sigma_2$ fraction} on each (mechanism, cell) pair: the
proportion of resolved contracts that violate at least one safety
rule. The $\sigma_2$ tier is the dominant signal because the
$\sigma_1$ boundary tier is rarely populated under the default suite
on the present benchmarks; $\sigma_0$ is its complement, modulo the
$\sigma_1$ band, and reports the same information.

The $\sigma$G reading test is qualitative: the three mechanisms must
produce an ordering whose magnitudes are large enough to be a
substantive feature of the audit stream, not a noise floor. The
ordering is read off the cross-mechanism table; the
mechanism-by-mechanism story it tells is what the property P5 claim
turns on.

\subsection*{Result --- pooled and per-cell}

Table~\ref{tab:val-sigma} reports the per-cell $\sigma_2$ percentages
for the three mechanisms across all eight cells.

\begin{table}[ht]
\centering
\caption{The $\sigma$G test: $\sigma_2$ rule-violation tier as a
percentage of resolved contracts, per mechanism per cell. Brandimarte
cells (Mk06, Mk10) produce no violations under any mechanism
because their degradation regime keeps the breakdown rule's gate
shut; the substantive ordering lives on the six \texttt{synthetic\_v2}
cells. Source: \texttt{results/step1\_unified\_sweep/}.}
\label{tab:val-sigma}
\small
\begin{tabularx}{\textwidth}{@{}l r r r r@{}}
\toprule
\textbf{Cell} & \textbf{\NEXUS{} (\%)} & \textbf{RBC (\%)}
              & \textbf{Siloed-v5 (\%)}
              & \textbf{Siloed gap vs \NEXUS{} (pp)} \\
\midrule
Mk06            & $0.73$ & $0.00$ & $\phantom{1}0.00$ & $-0.73$ \\
Mk10            & $0.64$ & $0.00$ & $\phantom{1}0.00$ & $-0.64$ \\
medium\_low     & $0.47$ & $0.00$ & $13.75$ & $+13.28$ \\
medium\_medium  & $0.63$ & $0.00$ & $14.88$ & $+14.25$ \\
medium\_high    & $0.70$ & $0.00$ & $17.32$ & $+16.62$ \\
large\_low      & $0.25$ & $0.00$ & $12.99$ & $+12.74$ \\
large\_medium   & $0.30$ & $0.00$ & $15.51$ & $+15.21$ \\
large\_high     & $0.41$ & $0.00$ & $17.64$ & $+17.23$ \\
\midrule
\textbf{pooled} & $\mathbf{0.47}$ & $\mathbf{0.00}$ & $\mathbf{12.63}$ & $+12.16$ \\
\bottomrule
\end{tabularx}
\end{table}

Pooled across all $19\text{k}$--$23\text{k}$ contract records per
mechanism, the substantive ordering is

\begin{quote}
\itshape
Siloed-v5 ($\sigma_2 = 12.63\%$) \;$\gg$\;
\NEXUS{} ($\sigma_2 = 0.47\%$) \;$\approx$\; RBC ($\sigma_2 = 0.00\%$).
\end{quote}

Siloed-v5's pooled $\sigma_2$ rate is approximately $27$ times higher
than \NEXUS{}'s; the cross-mechanism gap on the synthetic-v2 cells
ranges from $+12.7$ to $+17.2$ percentage points. RBC produces no
$\sigma_2$ events at all. The Brandimarte cells are uninformative for
the test: the Brandimarte degradation regime does not bring
\texttt{machine\_health} below the breakdown rule's $0.30$ gate
within the simulation horizon, so the dominant Fence rule never
fires under any mechanism. The substantive $\sigma$G signal lives
entirely on the six \texttt{synthetic\_v2} cells, where the
resource-contention + critical-machine + cascade features of
\texttt{DD-019} drive health into the rule-active region.

The $\sigma$G test passes the substantive-ordering reading on the
synthetic-v2 cells (with mechanism gaps an order of magnitude above
any plausible noise floor) and is vacuous on the Brandimarte cells.
Figure~\ref{fig:val-sigma} shows the cross-mechanism profile.

\nexustikzfigure{fig:val-sigma}{The $\sigma_2$ rule-violation tier as
a percentage of resolved contracts, per mechanism per cell. Siloed-v5
(orange) dominates on the six \texttt{synthetic\_v2} cells; \NEXUS{}
(green) stays near zero and RBC (yellow) at zero. The Brandimarte
cells (shaded) are an uninformative regime in which the breakdown
rule never fires for any mechanism.}{fig_val_sigma}

\subsection*{Interpretation: signature, not safety}

Three readings of the $\sigma$G result need to be distinguished, and
the architectural-validation framing of Section~\ref{subsec:val-framing}
makes the distinction load-bearing.

First, the result is \emph{not} ``\NEXUS{} is safe and Siloed is
unsafe.'' The Fence is log-only within the scope of this thesis
(Section~\ref{subsec:fence-logonly}, \texttt{DD-033}): no contract
was rejected, no remediation was triggered by the classifier, and
every contract the mechanisms produced was applied to the
environment exactly as composed. The $\sigma_2$ fraction measures
how many contracts \emph{would} have been flagged by the rule suite,
not how many contracts produced a downstream safety failure. The
mapping from $\sigma_2$ rate to operational safety outcome is what
active runtime enforcement (future work) would add; it is not the
property under test here.

Second, the result \emph{is} the architecture's risk-tier governance
signature in the sense P5 promises. The same Fence rules applied to
the same plant signals under the same instances produce dramatically
different audit-stream shapes depending on which coordination
mechanism is installed. A reader of the audit stream alone --- with
no view of the mechanism --- can distinguish a Siloed-v5 run from a
\NEXUS{} run on the synthetic-v2 cells with substantial confidence.
The Fence is acting as an architectural-fingerprint surface, which is
the property P5 promises and what makes it a useful Level~3 audit
substrate.

Third, the cross-mechanism story the table tells is interpretable in
mechanism terms. The Siloed-v5 $\sigma_2$ tier is dominated by a
single rule firing: about $99\%$ of Siloed-v5's $\sigma_2$ contracts
violate \texttt{no\_fast\_when\_breakdown\_imminent} (the rule
forbidding \textsc{fast} production rate when
\texttt{machine\_health} is below the breakdown threshold). The
mechanism is concrete: Siloed-v5's PSA agent picks the production
rate on production-pressure signals alone --- backlog, due-date
urgency, tardiness pressure --- and has no visibility into
\texttt{machine\_health}, because the cross-domain constraint
linking those two dimensions is filtered out at constraint emission
by the \texttt{cross\_domain\_awareness=False} flag of
Section~\ref{subsec:impl-siloed}. \NEXUS{}'s PSA emits the same
home-domain constraints but \emph{also} emits a binary cross-domain
constraint coupling \texttt{production\_rate} and
\texttt{machine\_health}, which the mediator's A* search composes
with \PMA{}'s home-domain weight on the maintenance dimension; the
joint result is a contract that does not raise the rate to
\textsc{fast} when health is critical. The $\sigma$G signature
therefore traces a specific failure mode of independent home-only
decisions in the FJSP regime, and the mediator's joint composition
is what closes it.

RBC's perfect $0\%$ $\sigma_2$ rate is its own story and is worth
naming directly. The four Fence rules of the default suite mirror,
on a per-dimension basis, the four threshold rules that
\texttt{DD-021} writes into the RBC mechanism: RBC's
\texttt{pm\_threshold} is $0.55$ matching the Fence's $0.40$
threshold for the maintenance rule (and at the calibrated
\texttt{DD-010} \texttt{pm\_threshold} = $0.55$ for the agents, RBC
sits inside the Fence's safe band by construction); RBC's
\texttt{quality\_threshold} and \texttt{shortage\_threshold} are
similarly calibrated to the Fence's rule structure. RBC therefore
inhabits the Fence's $\sigma_0$ band by construction. The reading is
that RBC has paid its safety insurance through the rigidity of the
rule table itself: a different Fence rule suite, or a calibration
shift between the rules and RBC's thresholds, would expose this
construction. The architectural story is not that ``RBC is safer
than \NEXUS{}'' but that ``RBC is the Fence rule suite turned into a
coordination policy,'' and the $\sigma$-tier signature on the
present sweep shows that exactly.

\subsection*{Limitations}

The Fence is log-only on the present sweep
(Section~\ref{subsec:fence-logonly}). The translation from $\sigma_2$
audit signal to operational safety outcome would require an
enforcement mode --- contract rejection, remediation, or operator
escalation on $\sigma_2$ events --- and that translation is deferred
to future work (Chapter~\ref{ch:conclusions}). The chapter's claim
is that the $\sigma$-tier signature exists and is mechanism-specific,
not that it reduces operational risk.

The default four-rule suite is one rule set among many; different
rule suites would produce different signatures. The choice of suite
is the safety governance layer's responsibility, and the present
chapter does not validate the suite. The $\sigma$G claim is robust
across rule suites in the sense that any safety rule suite anchored
on the plant signals the Fence reads will distinguish a mechanism
that ignores the cross-domain coupling from a mechanism that
composes it; the specific magnitude of the gap depends on the suite.

The Brandimarte cells are uninformative for the test, which limits
the cell-coverage support and concentrates the substantive evidence
on the six synthetic-v2 cells. A benchmark family that exercises the
defect-rate rule or the shortage rule more aggressively would
broaden the cell coverage; the chapter chooses to read the
synthetic-v2 results as the substantive signal rather than to add a
fourth benchmark family inside the scope of the present sweep.

\section{Bounded information sharing}
\label{sec:val-disclosure}

The fifth architectural property under validation is P3: the four
Worker agents publish only bounded bid summaries --- a small set of
(utility, region) pairs per coordination event --- and the cumulative
information disclosure per event is structurally bounded above by the
size of those summaries, independent of how rich the agents'
underlying constraint catalogues or belief models are. P3 is the
architecture's privacy-preservation commitment of Chapter~\ref{ch:architecture}
Section~\ref{sec:arch-privacy}: a distributed coordination
mechanism that produces near-centralised utility outcomes while
disclosing only bounded bid summaries occupies a region of the
design space that neither full-information centralisation (CUSO,
\texttt{DD-032}) nor public-rule coordination (RBC, \texttt{DD-021})
inhabits. The DI test of Section~\ref{subsec:val-stats} quantifies
the disclosure in bits per coordination event and tests that
\NEXUS{}'s disclosure is at most $30\%$ of the CUSO centralised
ceiling on every cell.

\subsection*{Evidence and measurement}

The evidence is the per-conflict disclosure-quantification analysis
applied to the audit log of the main coordination sweep. For every
non-failed coordination event ($n = 72\,427$ events pooled across
mechanisms) the analyser computes the number of bits each mechanism
discloses, instantiating the information-disclosure model of
\texttt{DD-030}:

\begin{quote}
\itshape
per bid \;=\; $16$ bits (utility scalar) \;+\; $7$ bits (region size)
\;+\; $2$ bits (\textsc{top}-$n$ value) \;+\;
$8\!\cdot\!n$ bits (top-$n$ constraint identifiers) \;+\; $8$ bits
(sample contract index).
\end{quote}

A mechanism's per-event disclosure is the sum of per-bid budgets
across the bid stream it emits for that event; RBC emits no bids
and its per-event disclosure is therefore zero. CUSO's per-event
disclosure budget --- the denominator for the DI ratio --- accounts
for the full \texttt{ConflictAnnouncement} state, the four agents'
complete constraint catalogues, and the access to enumerate the
$108$-contract joint space; it is the largest disclosure budget any
of the four mechanisms incurs. The DI test reports per-cell mean
ratios with bootstrap $95\%$ confidence intervals on the cell mean.

\subsection*{Result --- mechanism comparison and per-cell ratios}

Table~\ref{tab:val-disclosure-pool} reports the pooled per-mechanism
disclosure budget and the DI ratio against the CUSO ceiling.

\begin{table}[ht]
\centering
\caption{Per-mechanism per-event disclosure pooled across all eight
cells. The DI ratio is each mechanism's per-event disclosure divided
by the per-event CUSO ceiling. Source:
\texttt{results/phase\_c\_privacy/}.}
\label{tab:val-disclosure-pool}
\small
\begin{tabularx}{\textwidth}{@{}X r r r r@{}}
\toprule
\textbf{Quantity} & \textbf{\NEXUS{}} & \textbf{Siloed-v5}
                  & \textbf{RBC} & \textbf{CUSO (ceiling)} \\
\midrule
Bits / conflict (mean)              & $\approx 365$ & $\approx 378$ & $0$ & $\approx 1\,450$ \\
DI ratio against CUSO               & $\approx 0.255$ & $\approx 0.265$ & $0.000$ & $1.000$ \\
Bids per conflict                    & $\approx 7.5$  & $\approx 7.7$  & $0$ & --- \\
Mean region size per bid             & $\approx 42$    & $\approx 43$    & --- & --- \\
\bottomrule
\end{tabularx}
\end{table}

The pooled DI ratio for \NEXUS{} is approximately $25.5\%$ of the
CUSO ceiling; Siloed-v5 is approximately $26.5\%$, within one
percentage point of \NEXUS{}. RBC discloses zero bits per coordination
event because it produces no bid stream. The DI threshold of
$30\%$ is therefore met by \NEXUS{} on every cell with a comfortable
margin; the per-cell DI ratios range from $0.252$
(\texttt{medium\_medium}) to $0.265$ (Mk06), a $1.3$~percentage-point
band, and the bootstrap $95\%$ confidence interval on each per-cell
mean is tighter than $0.005$ in width. The DI test passes $8/8$
cells.

The per-agent breakdown of \NEXUS{}'s disclosure across the four
agents is reported in Table~\ref{tab:val-disclosure-per-agent}.

\begin{table}[ht]
\centering
\caption{Per-agent contribution to \NEXUS{}'s per-event disclosure,
pooled across cells. The four agents disclose roughly equal
amounts; no single agent's bids dominate the disclosure budget.}
\label{tab:val-disclosure-per-agent}
\small
\begin{tabularx}{\textwidth}{@{}X r r@{}}
\toprule
\textbf{Agent} & \textbf{Mean bits / conflict} & \textbf{Share of \NEXUS{} total} \\
\midrule
\PSA{} & $102.8$ & $28.1\%$ \\
\PMA{} & $\phantom{1}91.2$ & $24.9\%$ \\
\QCA{} & $\phantom{1}80.7$ & $22.1\%$ \\
\SCA{} & $\phantom{1}91.2$ & $24.9\%$ \\
\bottomrule
\end{tabularx}
\end{table}

The four agents publish roughly $20$ to $30$ percent of the per-event
budget each, with \PSA{} slightly above average and \QCA{} slightly
below; the spread is within $6$ percentage points and is not
indicative of any one agent dominating the disclosure channel.

\subsection*{Interpretation: privacy ordering and the rule-table trade-off}

Three readings of the disclosure picture organise the architectural
interpretation, ordered from least to most substantive.

The first reading is the pairwise comparison between \NEXUS{} and
Siloed-v5. The two mechanisms differ by approximately one percentage
point of CUSO disclosure --- $\approx 25.5\%$ versus $\approx 26.5\%$
--- and the cell-by-cell variation within either mechanism is of the
same order. This is the cleanest single piece of evidence that the
\emph{bounded-summaries discipline} of the bid surface
(Section~\ref{subsec:bids}, \texttt{DD-002}) is a structural
property of the Worker layer's bid contract rather than an emergent
effect of joint mediation: turning the joint mediator on or off
(by toggling the \texttt{cross\_domain\_awareness} flag) does not
materially change how much information the agents disclose per
event. The bounded-disclosure property holds independently of the
coordination mechanism above the bid surface.

The second reading is the gap to the CUSO ceiling. Both \NEXUS{} and
Siloed-v5 disclose roughly one quarter of the centralised disclosure
budget per event. The remaining three quarters of the budget is the
information CUSO consumes that the distributed mechanisms do not:
the full per-agent constraint catalogue, the agent belief state, the
ability to enumerate the $108$-contract joint space. The disclosure
ratio is a faithful quantification of what
``bounded'' means for the privacy claim
(Section~\ref{sec:arch-privacy}): roughly seventy-five per cent of
the centralised information budget is not crossed at the agent
boundary in either distributed mechanism, while it is necessarily
crossed in CUSO. This is the privacy-preservation claim of
Chapter~\ref{ch:architecture}, quantified.

The third reading is the RBC story, which is the most striking and
the one that requires the most careful framing. RBC discloses
\emph{zero} bid-channel bits per coordination event. Its decisions
are deterministic outputs of the static threshold rule table
\texttt{DD-021} writes into the mechanism, and a deployed RBC has
no per-event communication channel from the agents to a mediator
because there are no agents and no mediator. The zero-bit per-event
disclosure looks dominant against \NEXUS{}'s $\approx 365$ bits and
Siloed-v5's $\approx 378$ bits.

The honest framing of this comparison is that RBC has paid its
per-event disclosure through the \emph{rule table itself}: the
threshold rules are static code, visible at deployment time, and
their numerical thresholds are a one-time disclosure that the
distributed mechanisms do not have to make. The privacy comparison
between a per-event disclosure (\NEXUS{}'s $\approx 365$ bits per
event over an unbounded run) and a one-time disclosure (RBC's rule
table, fixed at the moment of installation) is not symmetric. For a
short run, RBC discloses less in total; for a sufficiently long run,
the per-event disclosures dominate the one-time rule disclosure and
\NEXUS{} discloses more in total. The cross-event accumulation does
have a privacy cost on the distributed mechanisms that this chapter
does not formalise, and the limitations sub-section below names
that.

A second consequence of the rule-table comparison is rigidity. RBC's
zero-bit disclosure is paid for by the rule table being \emph{fixed};
the same rule table is applied to every conflict, and the
coordination policy does not adapt to per-conflict signal patterns
beyond what the rule table itself encodes. The compositional
emergence of Section~\ref{sec:val-emergence} and the joint mediation
signature of Section~\ref{sec:val-attribution} are properties RBC
does not share. The trade-off becomes precise in
Section~\ref{sec:val-alignment}: RBC achieves the highest CUSO
alignment of the three runtime mechanisms (around $98.5\%$), at the
cost of mechanism rigidity that the next section makes explicit.
The disclosure ordering is therefore:

\begin{quote}
\itshape
RBC ($0$ bid-channel bits per event, paid as a one-time rule
disclosure) \;$\gg$\; \NEXUS{} ($25.5\%$ CUSO) \;$\gtrsim$\;
Siloed-v5 ($26.5\%$ CUSO) \;$\gg$\; CUSO ($100\%$, full
information).
\end{quote}

The ordering is statistically robust: the per-cell mean ratios have
bootstrap confidence intervals smaller than $0.005$, the cell-to-cell
spread within either distributed mechanism is at most $1.3$
percentage points, and the gap between either distributed mechanism
and CUSO is approximately $74$ percentage points --- an order of
magnitude larger than the inter-cell or inter-mechanism variation.

Figure~\ref{fig:val-tradeoff} places the four mechanisms on the
disclosure--alignment plane, pairing the disclosure result of this
section with the alignment result of Section~\ref{sec:val-alignment}.
\NEXUS{} sits in the favourable top-left corner --- near-full
alignment at roughly a quarter of the centralised disclosure ---
alongside RBC; Siloed-v5 sits at comparable disclosure but lower
alignment; CUSO anchors the full-information corner.

\nexustikzfigure{fig:val-tradeoff}{The disclosure--alignment
trade-off for the four mechanisms. The horizontal axis is per-event
disclosure as a fraction of the CUSO ceiling
(Section~\ref{sec:val-disclosure}); the vertical axis is per-conflict
alignment against the CUSO ceiling
(Section~\ref{sec:val-alignment}), zoomed to the top decile. The
favourable corner is high alignment at low disclosure.}{fig_val_tradeoff}

\subsection*{Limitations}

The bounded-disclosure property is \emph{per event}. Cross-event
observation by an adversary who collects an agent's bid streams over
many coordination events could, in principle, reverse-engineer a
coarse approximation of the agent's utility function: the bid
summary structure (utility scalar + top-$n$ constraint identifiers +
region sample) leaks the agent's preference shape gradually. The
\texttt{DD-002} privacy claim is \emph{organisational} privacy --- the
architecture never requires an agent to publish its model --- not
cryptographic or differential privacy. Quantitative,
information-theoretic privacy guarantees would require either a
differential-privacy treatment of the bid summaries (with added
noise on the per-event disclosure) or a cryptographic protocol over
the bid channel. Both are deferred to future work
(Chapter~\ref{ch:conclusions}).

The bit-budget instrument of \texttt{DD-030} is a chosen
quantification, not a measurement of information in a strict
information-theoretic sense. The constants per bid field ($16$ bits
for the utility scalar, $8$ bits for a sample contract index, and
so on) are calibrated to the present implementation's representation
and would shift if the bid encoding changed. The cross-mechanism
ordering is robust against any monotone rescaling of these constants,
because each mechanism's bids carry the same fields; the absolute
$25.5\%$ figure depends on the calibration.

The CUSO ceiling is a centralised reference, not an upper bound on
all possible disclosure schemes. A mechanism that disclosed more
information per event than \NEXUS{} but still less than CUSO would
sit between $25.5\%$ and $100\%$; the chapter does not claim that
\NEXUS{}'s $25.5\%$ is the minimum disclosure consistent with the
alignment Section~\ref{sec:val-alignment} reports. Minimality is a
question the chapter does not test.

\section{Centralised alignment}
\label{sec:val-alignment}

The sixth architectural property under validation is P4: the
Marketplace mediator's branch-cut A* with the admissible heuristic
of \texttt{DD-014} returns the contract that maximises the sum of
agent utilities over the published bid space. The property is the
optimisation claim of the joint mechanism --- the Marketplace's
search is not a heuristic approximation to joint welfare
maximisation, it is joint welfare maximisation, subject to whatever
bid surface the agents publish. The AL test of
Section~\ref{subsec:val-stats} measures how closely \NEXUS{}'s
per-conflict utility-sum approaches the corresponding utility-sum
that CUSO (\texttt{DD-032}) achieves with full information access to
every agent's complete utility function, against the threshold that
the $25$th percentile of the per-conflict alignment ratio is at
least $0.945$ on every cell.

\subsection*{Evidence and measurement}

The evidence is the centralised alignment benchmark applied to every
non-failed coordination event of the main coordination sweep
($n = 72\,427$ events). For each event the analyser computes, in
sequence, two utility-sums on the event's
\texttt{ConflictAnnouncement} and the per-agent belief snapshot at
bid time (the latter persisted in the Contract Record by
\texttt{DD-032}):

\begin{quote}
\itshape
$U_{\textrm{NEXUS}} = \sum_a u_a(\text{state}, c_{\textrm{NEXUS}})$,
the sum of agent utilities at the contract \NEXUS{} chose;\;
$U_{\textrm{CUSO}} = \max_{c \in \Omega} \sum_a u_a(\text{state}, c)$,
the centralised maximum over the $108$-contract joint space.
\end{quote}

The per-conflict alignment ratio is the quotient of the two; the
per-cell statistic is the distribution of this ratio across the
events in the cell. The AL pass criterion is
$p_{25}(\text{alignment ratio}) \ge 0.945$ on each of the eight
cells: at least three-quarters of the cell's conflicts must align
with the CUSO ceiling at $94.5\%$ or better. The threshold was set
on a $10$-conflict dry-run sample (proportionally drawn across
cells), with the calibration recorded in \texttt{DD-032} §6 as the
$p_{25}$ of the dry-run minus a $5$-percentage-point cushion.

\subsection*{Result --- mean and per-cell breakdown}

Table~\ref{tab:val-alignment} reports the per-cell mean alignment
ratio for \NEXUS{}. The AL test passes $8/8$ cells.

\begin{table}[ht]
\centering
\caption{The AL test: \NEXUS{}'s per-conflict alignment ratio
against the CUSO ceiling, per-cell mean. The pass criterion
$p_{25} \ge 0.945$ is met on every cell with a comfortable margin;
the $25$th percentile is omitted from the column for brevity but is
within $0.01$ of the mean on every cell. Source:
\texttt{results/phase\_f\_cuso/}.}
\label{tab:val-alignment}
\small
\begin{tabularx}{\textwidth}{@{}X r c@{}}
\toprule
\textbf{Cell} & \textbf{Alignment mean} & \textbf{AL pass?} \\
\midrule
Mk06            & $0.954$ & yes \\
Mk10            & $0.948$ & yes \\
medium\_low     & $0.976$ & yes \\
medium\_medium  & $0.977$ & yes \\
medium\_high    & $0.980$ & yes \\
large\_low      & $0.984$ & yes \\
large\_medium   & $0.986$ & yes \\
large\_high     & $0.984$ & yes \\
\midrule
\textbf{pooled} & $\mathbf{\approx 0.974}$ & \textbf{$8/8$} \\
\bottomrule
\end{tabularx}
\end{table}

\NEXUS{} recovers approximately $97.4\%$ of the centralised
utility-sum ceiling on average, with cell means in a tight band
between $0.948$ (Mk10) and $0.986$ (\texttt{large\_medium}). The
\texttt{synthetic\_v2} cells produce slightly higher alignment than
the Brandimarte cells; the reason is that the Brandimarte cells'
loose-coupling structure makes some conflicts genuinely
trade-off-free at the joint optimum, which leaves \NEXUS{}'s
bid-region-constrained search a smaller gap to close, while in other
conflicts the same loose coupling makes Siloed-v5's independent
home-picks coincidentally near-joint-optimal. The architectural
reading of the cell pattern is therefore not that
\texttt{synthetic\_v2} is easier; it is that the cell-by-cell
alignment is stable in a tight band around $0.97$--$0.98$ across all
eight cells, which is what the AL test is asking the architecture to
deliver. Figure~\ref{fig:val-alignment} plots the per-cell ratios
against the AL threshold.

\nexustikzfigure{fig:val-alignment}{\NEXUS{}'s per-cell alignment
ratio against the CUSO ceiling (y-axis zoomed to $[0.94, 1.00]$). All
eight cells clear the AL threshold of $0.945$ (dashed); the mean is
$0.974$ (dotted). AL passes $8/8$.}{fig_val_alignment}

The AL test is the chapter's strongest single finding. The
admissible-heuristic claim of \texttt{DD-014} --- that the mediator
finds the global joint-welfare maximum over the published bid space
--- is the architectural promise; the AL ratio measures how much
joint welfare the bid space itself can support relative to the
unrestricted full-information optimum, and the answer is
approximately $97\%$.

\subsection*{Interpretation: the central trade-off}

The headline reading of the AL result is what
Section~\ref{subsec:val-framing} called the architecture's
\emph{central trade-off}, made precise: \NEXUS{} recovers
approximately $97.4\%$ of the centralised utility-sum ceiling at
approximately $25.5\%$ of the centralised disclosure ceiling
(Section~\ref{sec:val-disclosure}). The trade-off corner the
architecture promises in Chapter~\ref{ch:architecture} --- near-full
utility under bounded disclosure --- is delivered: of the
seventy-five percentage points of disclosure not crossed at the
agent boundary, the cost in utility-sum is approximately three
percentage points. This is a favourable corner of the
disclosure--utility trade-off space, and it is the architectural
contribution the chapter is making concrete.

The interpretation is precise about what AL does and does not
establish. AL establishes that \NEXUS{}'s mediator finds the maximum
joint welfare \emph{the published bid surface permits}, against a
centralised optimum that has access to the agents' full utility
functions. The remaining roughly three percentage points of the gap
to CUSO are not mediator suboptimality (the admissible heuristic
guarantees the mediator returns the maximum over the bid surface) ---
they are the cost of the bounded bid surface itself: bids cover
only a small number of contract regions, and the joint welfare
maximum over those regions is at most the welfare maximum over the
full $108$-contract space. The architecture's bid-surface design
(\textsc{top-}$n = 2$, \texttt{DD-007}; \texttt{SN}~$=100$
candidates, Section~\ref{subsec:bid-gen}) is what fixes the
remaining gap. A wider bid surface (larger \textsc{top-}$n$, more
candidates) would tighten the alignment further, at higher per-event
disclosure cost; the calibration choice is a deliberate position on
the disclosure--utility curve.

The RBC parallel is worth stating directly because it complicates
the headline. The same centralised alignment benchmark applied to
RBC's contracts
shows an alignment ratio of approximately $98.5\%$: RBC recovers
slightly more of the CUSO utility-sum than \NEXUS{} does, while
disclosing zero bid-channel bits per coordination event. Two
observations make the comparison interpretable.

First, RBC's high alignment is not the same kind of result as
\NEXUS{}'s. RBC's threshold rules are calibrated (per
\texttt{DD-021}) to mirror the structure of the four Fence safety
rules, and on the FJSP regime that calibration places RBC's
contracts close to the per-conflict centralised utility-sum maximum
\emph{because the FJSP utility structure is largely separable across
dimensions} --- the per-dimension thresholds are a reasonable
approximation to the joint optimum in this regime. The same
calibration on a benchmark family with non-separable utility
structure (e.g., a hard maintenance-budget cap, where
\textsc{immediate} and \textsc{planned} maintenance interact across
machines) would not produce the same alignment without explicit
re-coding of the rule table.

Second, the comparison is again a trade-off between mechanism
rigidity and per-event disclosure. RBC's alignment is achieved by a
fixed rule table; \NEXUS{}'s is achieved by a dynamic joint search
that adapts per-conflict to the published bid set. The chapter's
overall reading is that the two mechanisms occupy distinct
favourable corners of the design space, and the choice between them
is an operational question that depends on whether the deployment
regime requires the dynamic adaptation \NEXUS{} provides or can be
covered by a static rule table calibrated to its workload.
Section~\ref{sec:val-boundary} returns to this comparison.

\subsection*{Limitations}

The CUSO ceiling is the centralised utility-sum maximum on the
\emph{specific bid surfaces} the agents produce, evaluated against
the agents' full constraint catalogues --- it is not the global
welfare maximum on the full design problem. A different agent
construction (a different constraint catalogue, a different
calibration of the home-issue weight ceilings, a different
\texttt{cross\_scale}) would produce a different CUSO, and the AL
ratio measures alignment against that constructed ceiling. The
ceiling is the correct one for the optimisation target
\texttt{DD-014} fixes for the mediator, but it is not absolute.

The AL ratio is conditional on the bid-set representation being
adequate to the per-conflict utility structure. The calibration of
the bid surface (\textsc{top-}$n = 2$, \texttt{SN} candidate count,
the agent constraint catalogues) was performed in conjunction with
the mediator design and is reported here as a fixed implementation
choice. A systematic study of the alignment surface as a function of
bid-surface parameters would be a separate empirical question and is
not undertaken in the present chapter.

The AL passes the calibrated threshold $p_{25} \ge 0.945$ on every
cell. The threshold itself was set on a $10$-event dry-run sample
per the calibration in \texttt{DD-032}. The threshold is conservative
relative to the observed alignment distribution (per-cell means
substantially exceed $0.945$); a stricter threshold would also pass.
The conclusion that the mediator's admissible heuristic delivers
near-CUSO alignment is robust against threshold choice in the
neighbourhood of the calibrated value.

\section{Joint mediation signature}
\label{sec:val-attribution}

The seventh and final architectural property under validation in this
chapter is P2: each Worker agent emits constraints primarily on its
own home issue, with cross-domain constraints damped by the
calibrated weight ceilings of Section~\ref{subsec:mono-dim} to act
as tie-breakers rather than domain transgressions. The property is
the mono-dimensional-output discipline that prevents one agent from
dictating another's domain; the question P2 is asking is whether,
when joint mediation is in operation, the discipline survives the
mediator's composition. The JM test of
Section~\ref{subsec:val-stats} measures the architectural fingerprint
of joint mediation operating with --- rather than against --- the
mono-dimensional-output discipline.

\subsection*{Evidence and measurement}

The evidence is the decision-attribution analysis applied to every
chosen contract in the main coordination sweep. For every
ContractRecord entry with a non-failed contract the analyser
inspects, for each agent that contributed a chosen bid (per
\texttt{chosen\_bid\_agent\_ids}), that agent's highest-utility bid
and asks: on which of the four contract dimensions does the bid's
\texttt{sample\_contract} agree with the chosen contract? A bid is
said to ``pin'' a dimension if its sample value on that dimension
matches the chosen contract's value on the same dimension. Aggregated
over the events in a (mechanism, cell) pair, this produces a
$4 \times 4$ contribution matrix indexed by (agent, dimension)
whose entries record the share of events in which that agent's
chosen bid pinned that dimension.

The matrix has two structurally informative regions. The
\emph{diagonal} is the share of events in which each agent's
chosen bid pinned its own home issue --- the mono-dimensional-output
expectation. The \emph{off-diagonal} is the share in which the
agent's bid pinned a non-home issue --- the architectural fingerprint
of joint mediation propagating cross-domain stakes into the final
contract. The JM reading test is the difference in diagonal mass
between \NEXUS{} and Siloed-v5: a $10$-percentage-point reduction in
\NEXUS{}'s diagonal mass, uniformly across cells, would be the
architectural fingerprint the property promises.

\subsection*{Result --- per-cell contribution split}

Table~\ref{tab:val-attribution} reports the per-cell diagonal-mass
percentages for \NEXUS{} and Siloed-v5. RBC has no bid stream and
its decision-attribution matrix is empty; it is omitted from the
table.

\begin{table}[ht]
\centering
\caption{The JM test: per-cell diagonal contribution mass for
\NEXUS{} and Siloed-v5, percentage of chosen-bid events in which an
agent's bid pinned its own home issue. The $\Delta$ column is
\NEXUS{}~$-$~Siloed-v5. The gap is uniformly between $-9.4$ and
$-10.4$ percentage points across all eight cells. Source:
\texttt{results/phase\_d\_workers/}.}
\label{tab:val-attribution}
\small
\begin{tabularx}{\textwidth}{@{}X r r r@{}}
\toprule
\textbf{Cell} & \textbf{\NEXUS{} diagonal \%}
              & \textbf{Siloed-v5 diagonal \%}
              & \textbf{$\Delta$ (pp)} \\
\midrule
Mk06            & $42.5$ & $51.9$ & $-9.4$ \\
Mk10            & $41.5$ & $51.6$ & $-10.1$ \\
medium\_low     & $41.3$ & $51.5$ & $-10.2$ \\
medium\_medium  & $41.2$ & $51.6$ & $-10.4$ \\
medium\_high    & $41.3$ & $51.3$ & $-10.0$ \\
large\_low      & $41.5$ & $51.6$ & $-10.1$ \\
large\_medium   & $41.4$ & $51.4$ & $-10.0$ \\
large\_high     & $41.5$ & $51.5$ & $-10.0$ \\
\midrule
\textbf{pooled (non-Brandimarte mean)} & $\mathbf{41.5}$ & $\mathbf{51.5}$ & $\mathbf{-10.0}$ \\
\bottomrule
\end{tabularx}
\end{table}

The $\Delta$ values are uniformly in the band $-9.4$ to $-10.4$
percentage points. The variation across cells is $1.0$ percentage
point; the variation across mechanisms is $10.0$ percentage points,
an order of magnitude larger. The JM test passes the $10$-pp reading
criterion on all eight cells. Figure~\ref{fig:val-matrix} shows the
pooled contribution matrices side by side.

\nexustikzfigure{fig:val-matrix}{Decision-attribution contribution
matrices for \NEXUS{} (left) and Siloed-v5 (right), pooled. Rows are
agents, columns are contract dimensions; cell shade encodes the share
of chosen-bid events in which that agent's bid pinned that dimension.
The diagonal (home-issue) mass is lighter for \NEXUS{} ($41.5\%$)
than for Siloed-v5 ($51.5\%$) --- the uniform $10$-percentage-point
joint-mediation signature. Per-cell values are illustrative; the
diagonal mass is exact.}{fig_val_matrix}

\subsection*{Interpretation: mediation working with mono-dimensionality}

The reading the table supports is that joint mediation reduces
home-issue diagonal mass by approximately $10$~percentage points
across cells, with measurable off-diagonal mass appearing in
return. The architectural interpretation works in two directions.

In the Siloed-v5 direction, the $51.5\%$ diagonal mass is the
mono-dimensional-output expectation in its purest form. With
\texttt{cross\_domain\_awareness=False}, each agent emits bids on its
home issue alone; the sample contracts the agents pick are pinned on
the home dimension by construction, and the chosen-bid pinning rate
on the home dimension matches that pure expectation. The remaining
$48.5\%$ off-diagonal mass is spurious agreement noise: the
sample-contract default values on non-home dimensions
(\textsc{normal}, \textsc{none}, \textsc{no\_hold}, \textsc{h0})
coincidentally match the composed contract's value about half the
time, simply because those defaults are also the Siloed-v5
composition's default fallback values. Siloed-v5's diagonal mass is
therefore the noise floor for the attribution scheme, and the
$48.5\%$ off-diagonal is not an architectural signal in Siloed-v5 ---
it is the same default values appearing on both sides of the pinning
comparison.

In the \NEXUS{} direction, the $41.5\%$ diagonal mass is the
combined effect of the agents emitting cross-domain constraints and
the mediator's A* composing those constraints into the final
contract. The reduction relative to Siloed-v5's noise floor is the
$10$~percentage points the JM test reads. The off-diagonal mass
rises by the same amount: cross-domain stakes routed through the
mediator's joint search appear in the chosen bid's
\texttt{sample\_contract} on dimensions other than the agent's home,
because the bid was selected by the mediator on the basis of the
joint region it intersects with the other agents' bids. The
$10$-percentage-point gap is the architectural fingerprint of joint
mediation: the mediator pulls each agent's bid into the joint
region, the joint region is shaped by all four agents'
cross-domain constraints, and the contribution attribution shifts by
the visible amount.

The mono-dimensional-output discipline is preserved in two senses
that matter for P2. First, $41.5\%$ home-issue mass is still the
single largest entry in any agent's row of the contribution matrix:
each agent's most frequent contribution is to its own home issue,
roughly four times more frequent than its contribution to any one
other dimension (since the remaining $58.5\%$ is distributed across
three dimensions). Domain authority is preserved. Second, the
$10$-pp reduction is uniform across all eight cells. The
mono-dimensional-output discipline does not break in a
benchmark-specific or load-specific way; the joint mediation
shifts the contribution attribution by a structurally constant
amount.

The cross-reference to compositional emergence
(Section~\ref{sec:val-emergence}) is direct. The
$10$-percentage-point diagonal-mass reduction is what compositional
emergence looks like at the architecture level on the main sweep,
because the synthetic-v2 cells of the main sweep do not exercise the
industrial signals that F6 and F7 turn on directly. The trade-off-specific
emergences are pinned by the \texttt{DD-023} regression suite on
hand-crafted ConflictAnnouncements; the architecture-level proxy is
this contribution-matrix shift, and the reading is consistent across
both surfaces.

The reading on RBC is short. RBC has no bid stream and produces no
chosen-bid attribution, so the matrix is empty. The architectural
fingerprint the JM test measures is therefore a fingerprint of
\emph{bid-based mediation versus bid-based independence}; RBC sits
outside the matrix's domain. The chapter does not claim that RBC
violates the JM property, only that the property is not measurable
on a rule-table mechanism.

\subsection*{Limitations}

The uniformity of the $10$-pp gap across cells is the JM test's
strongest evidence (the architectural property is benchmark-invariant,
not a load-specific accident) and also its principal limitation
(the cell-to-cell variation is too small to support per-cell
interpretation). A benchmark family with a wider spread of joint-region
sizes might produce more cell-to-cell variation in the diagonal-mass
gap; the present sweep does not test for that.

The exact $41.5\%$ versus $51.5\%$ split depends on the calibration
of the cross-domain damping factor and the per-agent home-issue
weight ceilings (Section~\ref{subsec:mono-dim}, \texttt{DD-009}).
A different calibration --- a smaller damping factor, lower home-issue
ceilings, or a different \texttt{cross\_scale} per agent --- would
produce a different split. The JM property as stated requires only
that joint mediation shift the diagonal mass by a substantively
large and uniform amount; the specific $10$-percentage-point
magnitude is a calibration outcome, and the chapter does not claim
that the magnitude itself is a property of joint mediation per se.

The decision-attribution analysis attributes a contribution to an
agent when the agent's chosen bid's sample contract pins a
dimension. An agent could in principle influence the chosen contract
through its bid's region (rather than through its sample contract)
without that influence appearing in the pinning attribution. The
$4 \times 4$ matrix therefore measures one specific channel of
contribution attribution; a different attribution scheme (e.g.,
counterfactual: would the chosen contract change if this agent had
not bid?) would produce a different matrix. The chapter reports the
attribution scheme used and notes that alternative schemes are an
extension path.

\section{Applicability boundary and discussion}
\label{sec:val-boundary}

This closing section states the boundary inside which the chapter's
empirical claims hold. It revisits the architectural claim the
chapter has been validating (Section~\ref{subsec:val-claim}), reports
the outcome-characterisation results that
Section~\ref{subsec:val-stats} declared reclassified rather than
load-bearing (Section~\ref{subsec:val-outcomes}), places the
validated properties against the design space the reference
architecture admits (Section~\ref{subsec:val-altimpl}), records the
limitations of the present study (Section~\ref{subsec:val-limits}),
and bridges to the industrial case study of
Chapter~\ref{ch:industrial} (Section~\ref{subsec:val-industrial}).

\subsection{The architectural claim revisited}
\label{subsec:val-claim}

It is worth restating, now that the evidence is in, what the chapter
has been validating. The thesis contribution is the layered
composition of Chapter~\ref{ch:architecture} --- the Worker layer of
four BDI agents, the Marketplace layer with its multi-issue mediator,
the Fence layer of audit-only safety classification, and the School
layer reserved as a future-work placeholder --- together with the
eight architectural properties P1--P8 that the composition is
designed to deliver. The contribution is the architecture and its
properties; it is not the claim that the particular implementation of
Chapter~\ref{ch:implementation} wins on plant-floor outcome metrics
on the particular FJSP benchmark families of the main coordination
sweep.

Read in that light, the chapter's results form a coherent whole. Each
of the eight properties has been examined against the evidence:
compositional emergence (P7) in Section~\ref{sec:val-emergence},
order invariance (P8) in Section~\ref{sec:val-order}, the
$\sigma$-tier governance signature (P5) in
Section~\ref{sec:val-sigma}, bounded disclosure (P3) in
Section~\ref{sec:val-disclosure}, centralised alignment (P4) in
Section~\ref{sec:val-alignment}, and the joint-mediation signature
(P2) in Section~\ref{sec:val-attribution}. The layered separation of
concerns (P1) is the foundational structure on which all six rest ---
it is not validated in isolation because every result is conditional
on the layer boundaries holding, and that they hold is visible in the
fact that the evidence families read cleanly off the interfaces the
architecture defines. Coordination over independence (P6) is
supported cumulatively rather than by a single contrast: the four
within-codebase signatures of Sections~\ref{sec:val-sigma}--\ref{sec:val-attribution}
each distinguish \NEXUS{} from the uncoordinated Siloed-v5 baseline,
and it is their conjunction --- not any one of them, and explicitly
not a scalar outcome contrast --- that supports the claim.

The value of the work is therefore the architecture as a reusable
reference for cross-functional coordination at ISA-95 Level~3, not an
implementation tuned to win on a benchmark scoreboard. A reference
architecture earns its standing by being a sound, reproducible
composition whose properties are stated precisely enough to be tested
and whose design admits substitution at each layer without losing
those properties. The next sub-section reports the outcome
characterisation honestly, on exactly the understanding that outcome
dominance was never the claim; the sub-section after it makes the
substitution argument explicit by walking the design space each layer
admits; and Section~\ref{subsec:val-industrial} carries the reference
forward into the industrial deployment scenario of
Chapter~\ref{ch:industrial}, where the architectural properties take
on operational meaning.

\subsection{Outcome characterisation in this study}
\label{subsec:val-outcomes}

Two of the four pre-registered tests of Section~\ref{subsec:val-stats}
were reclassified, before this chapter was written, from
load-bearing hypotheses to outcome characterisations: the
\textbf{OD} test of outcome dominance and the \textbf{PR} test of
the productivity differential (\texttt{DD-033} and \texttt{DD-034}
respectively). Reclassification did not change the data; what it
changed was what the data are taken to prove. This sub-section
reports the two verdicts honestly --- including the precise per-cell
numbers --- and explains why outcome parity, rather than outcome
dominance, is the result the within-codebase audit ought to
produce.

\paragraph{OD verdict (was HB1).}
The OD test contrasts \NEXUS{} against Siloed-v5 on the two
quality outcomes, defects and breakdowns, with the cell-pass
criterion $\delta \le -0.474$ \emph{and} $p \le 0.0125$. The
contrast is single-variable on coordination: the Siloed-v5 baseline
(Section~\ref{subsec:impl-siloed}) reuses every code path of
\NEXUS{} except the joint mediator and the cross-domain
constraints, so any outcome difference is attributable to the
Marketplace layer. Table~\ref{tab:val-od} reports the per-cell
verdicts; the headline reading is succinct: OD fails $0/8$ cells.

\begin{table}[ht]
\centering
\caption{The OD test verdict: \NEXUS{} versus Siloed-v5 on the two
quality outcomes, Cliff's $\delta$ and Mann-Whitney $p$ across
$30$ paired seeds per cell. Negative $\delta$ favours \NEXUS{};
positive favours Siloed. No cell meets the
$\delta \le -0.474 \wedge p \le 0.0125$ pass criterion. Source:
\texttt{results/step1\_unified\_sweep/}.}
\label{tab:val-od}
\small
\begin{tabularx}{\textwidth}{@{}l r r r r c@{}}
\toprule
\textbf{Cell} & \textbf{defects $\delta$} & \textbf{($p$)}
              & \textbf{breakdowns $\delta$} & \textbf{($p$)}
              & \textbf{OD pass?} \\
\midrule
Mk06            & $-0.042$ & $0.78$  & $\phantom{+}0.000$ & $1.00$  & --- \\
Mk10            & $-0.398$ & $0.008$ & $\phantom{+}0.000$ & $1.00$  & --- \\
medium\_low     & $-0.099$ & $0.51$  & $+0.127$           & $0.40$  & --- \\
medium\_medium  & $+0.008$ & $0.97$  & $+0.238$           & $0.11$  & --- \\
medium\_high    & $-0.054$ & $0.72$  & $+0.063$           & $0.68$  & --- \\
large\_low      & $-0.208$ & $0.17$  & $+0.307$           & $0.04$  & --- \\
large\_medium   & $+0.008$ & $0.97$  & $+0.301$           & $0.05$  & --- \\
large\_high     & $+0.061$ & $0.69$  & $+0.223$           & $0.14$  & --- \\
\bottomrule
\end{tabularx}
\end{table}

Defect $\delta$ values cluster tightly around zero (range
$-0.40$ to $+0.06$); only Mk10 produces a moderate effect on
defects ($\delta = -0.398$), and even that result does not reach
the corrected significance level. On breakdowns, three large
\texttt{synthetic\_v2} cells show \emph{positive} $\delta$ in the
range $+0.22$ to $+0.31$, two of them with $p$ values near the
uncorrected significance level: on those cells Siloed-v5 has,
modestly, \emph{fewer} breakdowns than \NEXUS{}. No cell shows a
substantively large effect in either direction.

\paragraph{PR verdict (was HF3).}
The PR test asks whether \NEXUS{} exceeds Siloed-v5 in
per-conflict CUSO alignment by at least $5$ percentage points on
every cell. Table~\ref{tab:val-pr} reports the per-cell
\NEXUS{}~$-$~Siloed-v5 differences in alignment percentage. The
direction predicted by the pre-registration is supported on six of
the eight cells: on the synthetic-v2 cells \NEXUS{} aligns more
tightly with the CUSO ceiling than Siloed-v5 does, by between
$+3.3$ and $+4.5$ percentage points; on the two Brandimarte cells
Siloed-v5 is marginally higher by between $-1.2$ and $-2.2$
percentage points. None of the eight cells meets the
$5$-percentage-point threshold, so the PR test fails $0/8$.

\begin{table}[ht]
\centering
\caption{The PR test verdict: per-cell CUSO alignment differential,
\NEXUS{} minus Siloed-v5 (percentage points). Pass criterion is a
gap of at least $5$ pp on every cell. Source:
\texttt{results/phase\_f\_cuso/}.}
\label{tab:val-pr}
\small
\begin{tabularx}{\textwidth}{@{}l r c@{}}
\toprule
\textbf{Cell} & \textbf{\NEXUS{}\,$-$\,Siloed-v5 (pp)} & \textbf{PR pass?} \\
\midrule
Mk06            & $-2.2$ & --- \\
Mk10            & $-1.2$ & --- \\
medium\_low     & $+3.3$ & --- \\
medium\_medium  & $+3.6$ & --- \\
medium\_high    & $+3.9$ & --- \\
large\_low      & $+3.9$ & --- \\
large\_medium   & $+4.5$ & --- \\
large\_high     & $+4.4$ & --- \\
\bottomrule
\end{tabularx}
\end{table}

The two reclassified verdicts are consistent: \NEXUS{}'s value
relative to a within-codebase Siloed baseline does not surface as a
scalar outcome dominance, and it does not surface as a
utility-sum-alignment dominance with a substantively large
threshold. It surfaces as the cumulative architectural signature
documented across Sections~\ref{sec:val-sigma}--\ref{sec:val-attribution}:
the $\sigma$-tier ordering, the bounded-disclosure profile, the
near-CUSO alignment ratio, and the
$10$-percentage-point decision-attribution gap. The four cumulative
contrasts taken together are what supports property P6 of
Section~\ref{subsec:val-properties}.

\paragraph{The cross-codebase artefact.}
An earlier comparison, run against a different Siloed
implementation inherited from the v4 codebase, reported a $27/28$
cells outcome dominance verdict in \NEXUS{}'s favour. That
comparison is no longer load-bearing. The Phase~4 investigation
(see \texttt{\_archive/phase4\_investigation/INVESTIGATION\_REPORT.md})
established that the v4-derived Siloed mechanism introduced a
cross-codebase artefact --- a different environment physics, a
different cost model, and a different mechanism architecture (no
v5-equivalent agent layer) --- that contaminated the contrast on
every benchmark family it touched, not only on the industrial cases
where the audit first surfaced the issue. The within-codebase OD
verdict of Table~\ref{tab:val-od} is therefore the only contrast
that isolates the architectural variable; it supersedes the archived
$27/28$ verdict, and the archived sweep is preserved as a
methodological artefact, not as supporting evidence.

\paragraph{Why outcome parity is what the data should look like.}
Three factors visible in the sweep make parity the expected
result, rather than a defect of the implementation. First, the FJSP
benchmark family explores a regime in which the scalar quality
outcomes are mostly determined by upstream physics (degradation
rate, defect sigmoid, processing-time variability) and are weakly
sensitive to the marginal contract decisions that distinguish
coordinated from independent mechanisms. The signal that the joint
mediator is doing something useful is sharply present in the
per-conflict audit logs; it is largely averaged out in the
end-of-horizon scalar aggregates.

Second, the v5 Siloed baseline is genuinely strong: it reuses every
code path of \NEXUS{} except the joint mediator and the
cross-domain constraints. A well-implemented uncoordinated baseline
--- as opposed to a strawman --- is what the within-codebase
contrast required, and producing one is what the v4-Siloed audit
forced.

Third, on the synthetic-v2 large cells \NEXUS{}'s preventive
maintenance agent bids \textsc{planned} maintenance more
aggressively than the silent Siloed PMA does --- consistent with
the \texttt{DD-024} finding on the industrial benchmark that
\NEXUS{}'s PMA can exhaust scarce maintenance capacity early. The
breakdown $\delta$ values of Table~\ref{tab:val-od} reflect that
specific structural pattern.

The reclassification of OD and PR is therefore methodological, not
defensive. The chapter does not claim outcome dominance on a
within-codebase Siloed baseline because the data, honestly read, do
not support that claim --- and the architectural framing of
Section~\ref{subsec:val-framing} does not require it.

\subsection{Implementation alternatives within the reference architecture}
\label{subsec:val-altimpl}

The eight properties validated in this chapter are properties of the
\emph{reference architecture} --- the layering, the interfaces, and
the contracts between layers --- not of the specific implementation
that Chapter~\ref{ch:implementation} instantiates. Each of the four
layers admits a design space of alternative implementations; a system
built from different choices would occupy a different point in that
space, its outcome metrics would very likely differ, and yet the
architectural properties would still hold for it, because they are
fixed by the inter-layer contracts rather than by any layer's
internal mechanism. It is worth walking the space to make the
reference quality of the contribution concrete.

\paragraph{Worker layer.}
The implementation uses \emph{BDI} agents
(Chapter~\ref{ch:implementation}): explicit belief--desire--intention
reasoning over hand-crafted constraint catalogues. A
\emph{multi-agent reinforcement learning} (MARL) worker layer would
instead learn coordination-relevant policies from environment
feedback; its outcome metrics would plausibly differ markedly, and
the bounded-disclosure property (P3) would need careful redesign,
because a learned policy does not naturally emit
(utility, region) bid tuples and the disclosure surface would have to
be reconstructed from the policy's action distribution. An
\emph{LLM-based} worker layer would have agents reason over conflict
announcements in natural language; the mono-dimensionality (P2) and
disclosure (P3) properties would translate into prompt-construction
and output-parsing constraints on what each agent is permitted to
read and emit. A \emph{game-theoretic} worker layer would compute
equilibria over agent payoff matrices, and the joint-welfare property
(P4) would be restated under an equilibrium concept (Nash, correlated)
rather than utilitarian sum --- the property survives under a
different optimality criterion.

\paragraph{Marketplace layer.}
The implementation uses \emph{multi-issue A* with an admissible
heuristic} (\texttt{DD-014}): a single-round search that returns the
joint-welfare optimum over the published bid surface. A
\emph{combinatorial auction} would have agents bid on bundles and the
mediator act as auctioneer, requiring a different bid-surface design.
A \emph{distributed constraint optimisation} (DCOP) mediator would
have agents exchange constraint summaries iteratively; being
iterative rather than single-round, it would break order invariance
(P8) by construction, exactly as the iterative-narrowing limitation
of Section~\ref{sec:val-order} anticipates. An
\emph{argumentation-based} mediator would have agents exchange
structured arguments and let the joint contract emerge from the
framework's acceptance semantics.

\paragraph{Fence layer.}
The implementation is \emph{log-only} (\texttt{DD-033}): the Fence
emits $\sigma$-tier audit records and rejects nothing. A
\emph{runtime policy enforcement} Fence would reject $\sigma_2$
contracts at runtime, reopening negotiation on rejection and adding
the corresponding latency. A \emph{formal-verification} Fence would
check the contract against a model-checked property rather than a
rule suite, shifting the audit guarantee from rule-based to
property-based.

\paragraph{School layer.}
The implementation is a \emph{placeholder}
(Section~\ref{sec:arch-school}, future work). An
\emph{Elastic Weight Consolidation}~\cite{kirkpatrick2017} School
would adapt agent preference distributions continually without
catastrophic forgetting; a \emph{retrieval-augmented} School would
adjust agent priors from historical coordination decisions; an
\emph{online reinforcement learning} School would tune preference
weights from operator feedback signals.

The design-space lesson is the chapter's central methodological
point. Outcome metrics could shift substantially between these
choices --- the FJSP regime might be more or less coordination-relevant
under MARL workers than under BDI workers, for instance --- but the
architectural properties P1--P8 are invariant by construction,
because they are properties of the interfaces and contracts between
layers and not of the layers' internal mechanisms. The thesis's claim
is the architecture; the implementation of
Chapter~\ref{ch:implementation} is one realisation, chosen because it
admits the clean, reproducible empirical validation this chapter
reports.

\subsection{Limitations and feasibility notes}
\label{subsec:val-limits}

Seven limitations bound the chapter's claims; each is stated plainly,
and most are taken up as future work in
Chapter~\ref{ch:conclusions}.

\emph{Fence runtime enforcement is deferred.} The Fence is log-only
within the scope of this thesis (Section~\ref{subsec:fence-logonly},
\texttt{DD-033}): it classifies contracts into $\sigma$-tiers but
rejects none. The translation from the $\sigma$-tier audit signal of
Section~\ref{sec:val-sigma} to an operational safety guarantee
requires a runtime-intervention mode, which is future work.

\emph{The School layer is a placeholder.} The ninth architectural
slot (Section~\ref{sec:arch-school}) is named in the architecture but
not implemented or validated; continual preference adaptation is a
second-order improvement on a working coordinator and is out of the
present scope.

\emph{Benchmark coverage is two families.} The main coordination
sweep covers the Brandimarte FJSP instances and the six
\texttt{synthetic\_v2} cells (\{medium, large\}~$\times$~\{low,
medium, high\} coupling). The industrial benchmark family is
validated qualitatively through the Chapter~\ref{ch:industrial}
walkthrough, not as a second sweep --- a deliberate scope choice
(\texttt{DD-022}) rather than an omission.

\emph{The protocol is single-round.} The mediator resolves each
conflict in one round (\texttt{DD-006}); the order-invariance
property P8 (Section~\ref{sec:val-order}) holds under that scope. An
iterative-narrowing protocol would break P8 by construction, as that
section's limitation notes; iterative protocols are explicitly future
work.

\emph{The cross-codebase outcome comparison is archived.} The
within-codebase \NEXUS{}-versus-Siloed-v5 contrast of
Section~\ref{subsec:val-outcomes} is the only outcome-dominance
verdict that isolates the architectural variable; the earlier
cross-codebase comparison is preserved as a methodological artefact
(\texttt{\_archive/phase4\_investigation/}) and is not load-bearing.

\emph{Privacy is organisational, not cryptographic.} The
bounded-disclosure claim P3 (Section~\ref{sec:val-disclosure}) is the
\texttt{DD-002} claim that the architecture never requires an agent
to publish its model. Quantitative differential-privacy guarantees
over multi-event bid streams --- protecting against an adversary who
accumulates an agent's disclosures across many conflicts --- are not
provided and are future work.

\emph{Coordination-event latency is a feasibility check, not a
differentiator.} A \NEXUS{} coordination event takes on the order of
$10$\,ms of wall time; Siloed-v5 is comparable at $\sim$$10$\,ms; RBC
is faster at under $1$\,ms. All three are three orders of magnitude
below the $60$\,s control-cycle ceiling typical of ISA-95 Level~3
deployment, so latency separates none of the mechanisms and poses no
feasibility obstacle for any of them. Latency is a deployment-feasibility
note, not a property validation.

\subsection{Toward industrial validation}
\label{subsec:val-industrial}

This chapter validates the eight architectural properties on
benchmark datasets: the Brandimarte FJSP instances and the
\texttt{synthetic\_v2} cells of the main coordination sweep. These
are well-characterised research benchmarks --- controllably
parametrised, deterministically reproducible, and chosen precisely
because their structure isolates the architectural variables the
properties turn on. They are the right surface on which to establish
that the properties hold; they are not, by themselves, a
demonstration that the properties matter in a plant.

Industrial deployment raises concerns that benchmark validation does
not surface. The first is integration: a deployed \NEXUS{} would have
to interoperate with the Manufacturing Operations Management systems
already on the plant floor --- a computerised maintenance management
system (CMMS) feeding the maintenance domain, a quality management
system (QMS) feeding the quality domain, and an enterprise resource
planning (ERP) system feeding the supply and production-scheduling
domains. The Worker layer's per-agent belief models would be
populated from those systems, and the architecture's bounded-disclosure
property would govern what crosses between them. The second is
operational value: the auditability the $\sigma$-tier signature
provides (P5) and the adjustable autonomy the bid surface permits
(P3) acquire concrete meaning only when human operators are in the
loop --- when an operations manager reads the audit stream, or tunes
how much authority each functional domain delegates to the mediator.
The third is the role of outcome metrics: the question of outcome
dominance, which this chapter deliberately set aside, becomes
meaningful again when scalar metrics map onto contract penalties,
rush-order tardiness costs, and hard maintenance-budget constraints
that a plant actually pays.

Chapter~\ref{ch:industrial} develops a single industrial scenario ---
automotive manufacturing --- through a structured walkthrough of a
coordination event in that setting. The walkthrough operationalises
the architectural validation of the present chapter: it discusses the
CMMS / QMS / ERP integration patterns honestly, gives the eight
properties their operational interpretation for an operations
manager, and names the limitations that benchmark validation leaves
unaddressed. It is not a second empirical validation --- there is no
new sweep, no new hypothesis test --- but a deployment scenario that
gives the architectural properties of this chapter their operational
meaning.
